\chapter{2005年重庆卷（文）}
\section{单选题}
\begin{problem}
圆 \((x + 2)^{2} + y^{2} = 5\) 关于原点 \((0,0)\) 对称的圆的方程为（\quad）

\choices
    {\((x - 2)^{2} + y^{2} = 5\)}
    {\(x^{2} + (y - 2)^{2} = 5\)}
    {\((x + 2)^{2} + (y + 2)^{2} = 5\)}
    {\(x^{2} + (y + 2)^{2} = 5\)}
\end{problem}

\begin{answer}
A
\end{answer}

\begin{solution}
原圆的圆心为 \((-2,0)\)，关于原点对称后圆心变为 \((2,0)\)，半径不变，仍为 \(\sqrt{5}\). 因此所求圆的方程为 \((x-2)^2+y^2=5\)，故选 A.
\end{solution}

\begin{problem}
\(\left(\cos \frac{\pi}{12} - \sin \frac{\pi}{12}\right)\left(\cos \frac{\pi}{12} + \sin \frac{\pi}{12}\right)=\)（\quad）

\choices
    {\(-\frac{\sqrt{3}}{2}\)}
    {\(-\frac{1}{2}\)}
    {\(\frac{1}{2}\)}
    {\(\frac{\sqrt{3}}{2}\)}
\end{problem}

\begin{answer}
D
\end{answer}

\begin{solution}
利用平方差公式，原式等于 \(\cos^2\frac{\pi}{12}-\sin^2\frac{\pi}{12}=\cos\frac{\pi}{6}=\frac{\sqrt{3}}{2}\)，故选 D.
\end{solution}

\begin{problem}
若函数 \( f(x) \) 是定义在 \( \R \) 上的偶函数，在 \((- \infty, 0]\) 上是减函数，且 \( f(2) = 0 \)，则使得 \( f(x) < 0 \) 的 \( x \) 的取值范围是（\quad）

\choices
    {\((-\infty, 2)\)}
    {\((2, +\infty)\)}
    {\((-\infty, -2) \cup (2, +\infty)\)}
    {\((-2, 2)\)}
\end{problem}

\begin{answer}
D
\end{answer}

\begin{solution}
因为 \(f\) 在 \((-\infty,0]\) 上严格递减且为偶函数，\(f(-2)=f(2)=0\). 当 \(x\in(-2,0]\) 时，由严格递减性得 \(f(x)<f(-2)=0\)；当 \(x\in[0,2)\) 时，\(-x\in(-2,0]\)，从而 \(f(x)=f(-x)<0\). 当 \(x<-2\) 时，\(f(x)>f(-2)=0\)，再由偶性知当 \(x>2\) 时也有 \(f(x)>0\)，而 \(x=\pm2\) 时函数值为零. 因此 \(f(x)<0\) 的取值范围为 \((-2,2)\)，故选 D.
\end{solution}

\begin{problem}
设向量 \(\bs{a} = (-1, 2)\)，\(\bs{b} = (2, -1)\)，则 \((\bs{a} \cdot \bs{b})(\bs{a} + \bs{b})\) 等于（\quad）

\choices
    {\((1, 1)\)}
    {\((-4, -4)\)}
    {\(-4\)}
    {\((-2, -2)\)}
\end{problem}

\begin{answer}
B
\end{answer}

\begin{solution}
先求数量积：\(\bs{a}\cdot\bs{b}=(-1)\times2+2\times(-1)=-4\)，又 \(\bs{a}+\bs{b}=(1,1)\). 所以 \((\bs{a}\cdot\bs{b})(\bs{a}+\bs{b})=-4(1,1)=(-4,-4)\)，故选 B.
\end{solution}

\begin{problem}
不等式组 \( \begin{cases}
|x-2|<2,\\
\log_{2}(x^{2}-1)>1
\end{cases} \) 的解集为（\quad）

\choices
    {\((0, \sqrt{3})\)}
    {\((\sqrt{3}, 2)\)}
    {\((\sqrt{3}, 4)\)}
    {\((2,4)\)}
\end{problem}

\begin{answer}
C
\end{answer}

\begin{solution}
由 \(|x-2|<2\) 得 \(0<x<4\). 对数不等式还要求 \(x^2-1>0\)，且由于底数 \(2>1\)，有 \(\log_2(x^2-1)>1\Longleftrightarrow x^2-1>2\)，即 \(x^2>3\). 与 \(0<x<4\) 交集为 \((\sqrt{3},4)\)，故选 C.
\end{solution}

\begin{problem}
已知 \(\alpha\)，\(\beta\) 均为锐角，若 \(p: \sin \alpha < \sin (\alpha + \beta)\)，\(q: \alpha + \beta < \frac{\pi}{2}\)，则 \(p\) 是 \(q\) 的（\quad）

\choices
    {充分而不必要条件}
    {必要而不充分条件}
    {充要条件}
    {既不充分也不必要条件}
\end{problem}

\begin{answer}
B
\end{answer}

\begin{solution}
由和差化积公式，\(\sin(\alpha+\beta)-\sin\alpha=2\cos\left(\alpha+\frac{\beta}{2}\right)\sin\frac{\beta}{2}\). 由于 \(\beta\) 为锐角，\(\sin\frac{\beta}{2}>0\)，所以命题 \(p\) 等价于 \(\alpha+\frac{\beta}{2}<\frac{\pi}{2}\). 若命题 \(q\) 成立，则 \(\alpha+\frac{\beta}{2}<\alpha+\beta<\frac{\pi}{2}\)，故 \(q\Rightarrow p\). 但取 \(\alpha=\frac{7\pi}{18}\)，\(\beta=\frac{\pi}{6}\)，有 \(\alpha+\frac{\beta}{2}=\frac{5\pi}{12}<\frac{\pi}{2}\)，命题 \(p\) 成立，而 \(\alpha+\beta=\frac{5\pi}{9}>\frac{\pi}{2}\)，命题 \(q\) 不成立. 因此 \(p\) 是 \(q\) 的必要而不充分条件，故选 B.
\end{solution}

\begin{problem}
对于不重合的两个平面 \(\alpha\) 与 \(\beta\)，给定下列条件：

\begin{circlelist}
\item 存在平面 \(\gamma\)，使得 \(\alpha\)，\(\beta\) 都垂直于 \(\gamma\)；
\item 存在平面 \(\gamma\)，使得 \(\alpha\)，\(\beta\) 都平行于 \(\gamma\)；
\item 存在直线 \(l \subset \alpha\)，直线 \(m \subset \beta\)，使得 \(l \parallel m\)；
\item 存在异面直线 \(l\)，\(m\)，使得 \(l \parallel \alpha\)，\(l \parallel \beta\)，\(m \parallel \alpha\)，\(m \parallel \beta\).
\end{circlelist}

其中，可以判定 \( \alpha \) 与 \( \beta \) 平行的条件有（\quad）

\choices
    {\(1\) 个}
    {\(2\) 个}
    {\(3\) 个}
    {\(4\) 个}
\end{problem}

\begin{answer}
B
\end{answer}

\begin{solution}
条件 1 不能判定两平面平行，因为两个相交平面也可以同时垂直于同一个平面. 例如取一个平面为水平面，两个不同的竖直平面即可满足条件 1 而相交.

条件 2 可以判定两平面平行：若 \(\alpha\) 和 \(\beta\) 都平行于 \(\gamma\)，则它们的法向量都与 \(\gamma\) 的法向量平行，从而 \(\alpha\) 与 \(\beta\) 的法向量平行. 又因两平面不重合，所以 \(\alpha\parallel\beta\).

条件 3 不能判定两平面平行. 若 \(\alpha\) 与 \(\beta\) 相交于直线 \(r\)，可分别在两个平面内取与 \(r\) 平行的直线 \(l\) 和 \(m\)，仍能满足 \(l\parallel m\).

条件 4 可以判定两平面平行. 若 \(\alpha\) 与 \(\beta\) 相交于直线 \(r\)，则既平行于 \(\alpha\) 又平行于 \(\beta\) 的直线，其方向只能平行于两平面的交线 \(r\)，所以题设中的 \(l\) 与 \(m\) 必相互平行，不可能异面. 这与条件 4 矛盾. 因此 \(\alpha\) 与 \(\beta\) 不相交；结合两平面不重合，只能有 \(\alpha\parallel\beta\).

所以满足条件 2 和条件 4，共有 2 个条件，故选 B.
\end{solution}

\begin{problem}
若 \((1 + 2x)^{n}\) 展开式中含 \(x^{3}\) 项的系数等于含 \(x\) 项的系数的\(8\) 倍，则 \(n\) 等于（\quad）

\choices
    {\(5\)}
    {\(7\)}
    {\(9\)}
    {\(11\)}
\end{problem}

\begin{answer}
A
\end{answer}

\begin{solution}
在 \((1+2x)^n\) 的二项展开式中，含 \(x\) 项的系数为 \(\mathrm C_n^1 2=2n\)，含 \(x^3\) 项的系数为 \(\mathrm C_n^3 2^3=8\mathrm C_n^3\). 由题意，\(8\mathrm C_n^3=8\times2n\)，所以 \(\mathrm C_n^3=2n\). 因含有 \(x^3\) 项，\(n\geqslant3\)，故可约去 \(n\)，得到 \(\frac{(n-1)(n-2)}{6}=2\)，即 \((n-1)(n-2)=12\). 解得 \(n=5\) 或 \(n=-2\)，舍去负值，故选 A.
\end{solution}

\begin{problem}
若动点 \((x, y)\) 在曲线 \(\frac{x^2}{4} + \frac{y^2}{b^2} = 1\)（\(b > 0\)）上变化，则 \(x^2 + 2y\) 的最大值为（\quad）

\choices
    {\(\begin{cases}
\frac{b^2}{4} +4, & 0 < b < 4,\\
2b, & b\geqslant 4.
\end{cases}\)}
    {\(\begin{cases}
\frac{b^2}{4} +4, & 0 < b < 2,\\
2b, & b\geqslant 2.
\end{cases}\)}
    {\(\frac{b^2}{4} + 4\)}
    {\(2b\)}
\end{problem}

\begin{answer}
A
\end{answer}

\begin{solution}
由椭圆方程得 \(x^2=4\left(1-\frac{y^2}{b^2}\right)\)，且 \(-b\leqslant y\leqslant b\). 因此 \(x^2+2y=4+2y-\frac{4y^2}{b^2}\). 这是关于 \(y\) 的开口向下的二次函数，其顶点横坐标为 \(y=\frac{b^2}{4}\).

当 \(0<b<4\) 时，\(\frac{b^2}{4}<b\)，顶点在区间 \([-b,b]\) 内，最大值为 \(4+\frac{b^2}{4}\). 当 \(b\geqslant4\) 时，顶点横坐标不小于右端点，二次函数在 \([-b,b]\) 上的最大值在 \(y=b\) 处取得，最大值为 \(2b\). 当 \(b=4\) 时两式均为 \(8\)，故选 A.
\end{solution}

\begin{problem}
有一塔形几何体由若干个正方体构成，构成方式如图所示. 上层正方体下底面的四个顶点是下层正方体上底面各连接中点，已知最底层正方体的棱长为 \(2\)，且该塔形的表面积（含最底层正方体的底面面积）超过 \(39\)，则该塔形中正方体的个数至少是（\quad）

\bitmapfigure[max width=0.56\linewidth]{img/2005/chongqing_liberal/q10_fig1.png}

\choices
    {\(4\)}
    {\(5\)}
    {\(6\)}
    {\(7\)}
\end{problem}

\begin{answer}
C
\end{answer}

\begin{solution}
设从下到上共有 \(n\) 个正方体，最底层棱长为 \(s_1=2\). 相邻两层中，上层正方体下底面是下层正方体上底面四边中点所围成的正方形，因此 \(s_{i+1}=\frac{s_i}{\sqrt{2}}\)，从而 \(s_i^2=4\left(\frac12\right)^{i-1}\).

每个正方体的四个侧面均露出. 底面总面积为 \(s_1^2\)，各层侧面面积总和为 \(4\sum_{i=1}^n s_i^2\). 除最底层底面外，水平方向露出的面积为 \(s_n^2+\sum_{i=1}^{n-1}(s_i^2-s_{i+1}^2)=s_1^2\)，所以塔形表面积为 \(T_n=2s_1^2+4\sum_{i=1}^n s_i^2=8+16\sum_{i=0}^{n-1}\left(\frac12\right)^i=40-\frac{32}{2^n}\).

当 \(n=5\) 时，\(T_5=39\)，不满足超过 \(39\)；当 \(n=6\) 时，\(T_6=39.5>39\). 由于 \(T_n\) 随 \(n\) 增大而增大，正方体的个数至少为 \(6\)，故选 C.
\end{solution}

\section{填空题}
\begin{problem}
若集合 \(A = \{x \in \R \mid x^2 - 4x + 3 < 0\}\)，集合 \(B = \{x \in \R \mid (x - 2)(x - 5) < 0\}\)，则 \(A \cap B =\fillinblank{}\).
\end{problem}

\begin{answer}
\((2,3)\)
\end{answer}

\begin{solution}
由 \(x^2-4x+3=(x-1)(x-3)\)，得 \(A=(1,3)\). 又由 \((x-2)(x-5)<0\)，得 \(B=(2,5)\). 因此 \(A\cap B=(2,3)\).
\end{solution}

\begin{problem}
曲线 \(y = x^3\) 在点 \((1,1)\) 处的切线与 \(x\) 轴，直线 \(x = 2\) 所围成的三角形的面积为\fillinblank{}.
\end{problem}

\begin{answer}
\(\frac{8}{3}\)
\end{answer}

\begin{solution}
曲线 \(y=x^3\) 的导数为 \(y'=3x^2\)，所以在 \((1,1)\) 处的切线斜率为 \(3\)，切线方程为 \(y-1=3(x-1)\)，即 \(y=3x-2\). 它与 \(x\) 轴交于 \(\left(\frac23,0\right)\)，与直线 \(x=2\) 交于 \((2,4)\)，而 \(x=2\) 与 \(x\) 轴交于 \((2,0)\). 三角形的底为 \(2-\frac23=\frac43\)，高为 \(4\)，所以面积为 \(\frac12\times\frac43\times4=\frac83\).
\end{solution}

\begin{problem}
已知 \(\alpha\)，\(\beta\) 均为锐角，且 \(\cos (\alpha + \beta) = \sin (\alpha - \beta)\)，则 \(\tan \alpha = \)\fillinblank{}.
\end{problem}

\begin{answer}
\(1\)
\end{answer}

\begin{solution}
由题意，\(\cos(\alpha+\beta)=\sin(\alpha-\beta)\). 展开两边，得 \(\cos\alpha\cos\beta-\sin\alpha\sin\beta=\sin\alpha\cos\beta-\cos\alpha\sin\beta\)，即 \((\cos\alpha-\sin\alpha)(\cos\beta+\sin\beta)=0\). 因为 \(\beta\) 为锐角，\(\cos\beta+\sin\beta>0\)，所以 \(\cos\alpha=\sin\alpha\). 又 \(\alpha\) 为锐角，故 \(\alpha=\frac\pi4\)，从而 \(\tan\alpha=1\).
\end{solution}

\begin{problem}
若 \(x^{2} + y^{2} = 4\)，则 \(x - y\) 的最大值是\fillinblank{}.
\end{problem}

\begin{answer}
\(2\sqrt{2}\)
\end{answer}

\begin{solution}
由柯西不等式，\(x-y=(x,y)\cdot(1,-1)\leqslant\sqrt{x^2+y^2}\sqrt{1^2+(-1)^2}=2\sqrt{2}\). 当 \((x,y)=\sqrt{2}(1,-1)\) 时满足 \(x^2+y^2=4\)，且等号成立. 因此 \(x-y\) 的最大值为 \(2\sqrt{2}\).
\end{solution}

\begin{problem}
若 \(10\) 把钥匙中只有 \(2\) 把能打开某锁，则从中任取 \(2\) 把能将该锁打开的概率为\fillinblank{}.
\end{problem}

\begin{answer}
\(\frac{17}{45}\)
\end{answer}

\begin{solution}
从 \(10\) 把钥匙中任取 \(2\) 把，共有 \(\mathrm C_{10}^2=45\) 种等可能取法. 不能打开锁的情形是两把都来自不能开锁的 \(8\) 把钥匙，共有 \(\mathrm C_8^2=28\) 种. 因此能打开锁的取法数为 \(45-28=17\)，所求概率为 \(\frac{17}{45}\).
\end{solution}

\begin{problem}
已知 \(A\left(-\frac{1}{2},0\right)\)，\(B\) 是圆 \(F:\left(x - \frac{1}{2}\right)^2 + y^2 = 4\)（\(F\) 为圆心）上一动点，线段 \(AB\) 的垂直平分线交 \(BF\) 于 \(P\)，则动点 \(P\) 的轨迹方程为\fillinblank{}.
\end{problem}

\begin{answer}
\(x^{2}+\frac{4}{3}y^{2}=1\)
\end{answer}

\begin{solution}
取直角坐标系，记圆心 \(F=\left(\frac12,0\right)\)，固定点 \(A=\left(-\frac12,0\right)\). 设 \(P\) 在直线 \(BF\) 上，并以从 \(F\) 指向 \(B\) 的单位向量为 \(\bs{e}\)，令 \(\overrightarrow{FP}=t\bs{e}\)，则 \(\overrightarrow{FB}=2\bs{e}\)，其中 \(t\) 暂为实数.

因为 \(P\) 在线段 \(AB\) 的垂直平分线上，所以 \(PA=PB\). 设 \(\bs{e}\) 与 \(x\) 轴正向的夹角为 \(\theta\)，则 \(P=F+t\bs{e}\)，从而 \(PA^2=1+2t\cos\theta+t^2\)，而 \(PB^2=(2-t)^2\). 由 \(PA^2=PB^2\) 得 \(t(2\cos\theta+4)=3\). 由于 \(2\cos\theta+4\geqslant2\)，所以 \(t>0\)，这也说明交点确实位于从 \(F\) 指向 \(B\) 的射线上. 又由 \(x-\frac12=t\cos\theta\)，所以 \(4t=4-2x\)，即 \(t=1-\frac{x}{2}\).

再由 \(t^2=PF^2=\left(x-\frac12\right)^2+y^2\)，得到 \(\left(x-\frac12\right)^2+y^2=\left(1-\frac{x}{2}\right)^2\)，化简为 \(x^2+\frac43y^2=1\). 反过来，椭圆上任一点都满足 \(t=1-\frac{x}{2}>0\) 及上述距离关系；取从 \(F\) 指向该点的单位向量为 \(\bs{e}\)，令 \(B=F+2\bs{e}\)，即可得到相应的圆上动点 \(B\) 和垂直平分线交点 \(P\). 故这就是动点 \(P\) 的轨迹方程.
\end{solution}

\section{解答题}
\begin{problem}
若函数 \( f(x) = \frac{1 + \cos 2x}{2\sin\left(\frac{\pi}{2} - x\right)} + \sin x + a^2 \sin \left(x + \frac{\pi}{4}\right) \) 的最大值为 \( \sqrt{2} + 3 \)，试确定常数 \( a \) 的值.

\begin{solution}
\(1+\cos2x=2\cos^2x\)，且定义域要求 \(\cos x\ne0\)，所以 \(\frac{1+\cos2x}{2\sin\left(\frac\pi2-x\right)}=\frac{2\cos^2x}{2\cos x}=\cos x\). 因此 \(f(x)=\cos x+\sin x+a^2\sin\left(x+\frac\pi4\right)=(\sqrt{2}+a^2)\sin\left(x+\frac\pi4\right)\). 由于 \(\sqrt{2}+a^2>0\)，且取 \(x=\frac\pi4+2k\pi\) 时原式有定义并使正弦值为 \(1\)，故函数最大值为 \(\sqrt{2}+a^2\). 由题设 \(\sqrt{2}+a^2=\sqrt{2}+3\)，得 \(a^2=3\)，所以 \(a=\pm\sqrt{3}\).
\end{solution}
\end{problem}

\begin{problem}
加工某种零件需经过三道工序，设第一、二、三道工序的合格率分别为 \(\frac{9}{10}\)，\(\frac{8}{9}\)，\(\frac{7}{8}\)，且各道工序互不影响.

\begin{enumerate}
\item 求该种零件的合格率；
\item 从该种零件中任取 \(3\) 件，求恰好取到一件合格品的概率和至少取到一件合格品的概率.
\end{enumerate}

\begin{solution}
\begin{enumerate}
\item 设一件零件合格的事件为 \(A\). 由于三道工序相互独立，该零件合格必须三道工序都合格，所以 \(P(A)=\frac9{10}\times\frac89\times\frac78=\frac7{10}=0.7\).
\item 任取的 \(3\) 件零件可看作独立重复试验，每件合格的概率为 \(p=0.7\)，不合格的概率为 \(1-p=0.3\). 恰好取到一件合格品的概率为 \(\mathrm C_3^1(0.7)(0.3)^2=0.189\). 至少取到一件合格品的概率用对立事件计算，为 \(1-(0.3)^3=0.973\).
\end{enumerate}
\end{solution}
\end{problem}

\begin{problem}
设函数 \( f(x) = 2x^{3} - 3(a + 1)x^{2} + 6ax + 8 \)，其中 \( a \in \R \).

\begin{enumerate}
\item 若 \( f(x) \) 在 \( x = 3 \) 处取得极值，求常数 \( a \) 的值；
\item 若 \( f(x) \) 在 \( (-\infty,0) \) 上为增函数，求 \( a \) 的取值范围.
\end{enumerate}

\begin{solution}
\begin{enumerate}
\item 函数可导，若在 \(x=3\) 处取得极值，则必须有 \(f'(3)=0\). 而 \(f'(x)=6x^2-6(a+1)x+6a=6(x-1)(x-a)\)，所以 \(0=f'(3)=12(3-a)\)，得 \(a=3\). 当 \(a=3\) 时，\(f'(x)=6(x-1)(x-3)\)，在 \(x=3\) 左侧为负、右侧为正，确实在 \(x=3\) 处取得极小值.
\item 由 \(f'(x)=6(x-1)(x-a)\)，当 \(x<0\) 时有 \(x-1<0\). 要使 \(f\) 在 \((-\infty,0)\) 上为增函数，需对一切 \(x<0\) 有 \(f'(x)>0\)，这等价于 \(x-a<0\) 对一切 \(x<0\) 成立，也就是 \(a\geqslant0\). 当 \(a\geqslant0\) 时确有 \(x-a<0\)，故 \(f'(x)>0\)；若 \(a<0\)，取 \(x\in(a,0)\)，则 \(f'(x)<0\). 故 \(a\in[0,+\infty)\).
\end{enumerate}
\end{solution}
\end{problem}

\begin{problem}
如图，在四棱锥 \(P - ABCD\) 中，底面 \(ABCD\) 为矩形，\(PD \perp\) 底面 \(ABCD\)，\(E\) 是 \(AB\) 上一点，\(PE \perp EC\). 已知 \(PD = \sqrt{2}\)，\(CD = 2\)，\(AE = \frac{1}{2}\)，求：

\begin{enumerate}
\item 异面直线 \(PD\) 与 \(EC\) 的距离；
\item 二面角 \(E - PC - D\) 的大小.
\end{enumerate}

\begin{solution}
\begin{enumerate}
\item 建立空间直角坐标系，令 \(A=(0,0,0)\)，\(B=(2,0,0)\)，\(D=(0,h,0)\)，\(C=(2,h,0)\)，\(P=(0,h,\sqrt2)\)，\(E=\left(\frac12,0,0\right)\). 由 \(PE\perp EC\) 得 \(\overrightarrow{EP}=\left(-\frac12,h,\sqrt2\right)\)，\(\overrightarrow{EC}=\left(\frac32,h,0\right)\)，从而 \(\overrightarrow{EP}\cdot\overrightarrow{EC}=h^2-\frac34=0\)，所以 \(h=\frac{\sqrt3}{2}\). 又 \(\overrightarrow{ED}=\left(-\frac12,h,0\right)\)，故 \(\overrightarrow{ED}\cdot\overrightarrow{EC}=h^2-\frac34=0\)，即 \(ED\perp EC\). 显然 \(ED\perp PD\)，所以 \(ED\) 是异面直线 \(PD\) 与 \(EC\) 的公垂线，且 \(ED=\sqrt{\frac14+h^2}=1\).
\item 取平面 \(EPC\) 和平面 \(DPC\) 的法向量. 令 \(\overrightarrow{PC}=(2,0,-\sqrt2)\)，\(\overrightarrow{PE}=\left(\frac12,-h,-\sqrt2\right)\)，\(\overrightarrow{PD}=(0,0,-\sqrt2)\)，则可取 \(\bs{n}_1=\overrightarrow{PC}\times\overrightarrow{PE}=\left(-\sqrt2h,\frac{3\sqrt2}{2},-2h\right)\)，\(\bs{n}_2=\overrightarrow{PC}\times\overrightarrow{PD}=(0,2\sqrt2,0)\). 代入 \(h=\frac{\sqrt3}{2}\)，有 \(|\bs{n}_1|=3\)，\(|\bs{n}_2|=2\sqrt2\)，且 \(|\bs{n}_1\cdot\bs{n}_2|=6\). 因此二面角的平面角 \(\theta\) 满足 \(\cos\theta=\frac{6}{3\times2\sqrt2}=\frac{\sqrt2}{2}\)，从而 \(\theta=\frac\pi4\).
\end{enumerate}
\end{solution}
\end{problem}

\begin{problem}
已知中心在原点的双曲线 \(C\) 的右焦点为 \((2,0)\)，右顶点为 \((\sqrt{3},0)\).

\begin{enumerate}
\item 求双曲线 \(C\) 的方程；
\item 若直线 \(l: y = kx + \sqrt{2}\) 与双曲线 \(C\) 恒有两个不同的交点 \(A\) 和 \(B\)，且 \(\overrightarrow{OA} \cdot \overrightarrow{OB} > 2\)（其中 \(O\) 为原点）. 求 \(k\) 的取值范围.
\end{enumerate}

\begin{solution}
\begin{enumerate}
\item 设双曲线的标准方程为 \(\frac{x^2}{a^2}-\frac{y^2}{b^2}=1\)，则右焦点横坐标为 \(c=2\)，右顶点横坐标为 \(a=\sqrt3\). 由 \(c^2=a^2+b^2\) 得 \(b^2=4-3=1\)，所以双曲线方程为 \(\frac{x^2}{3}-y^2=1\).
\item 将 \(y=kx+\sqrt2\) 代入双曲线方程，得 \(\left(\frac13-k^2\right)x^2-2\sqrt2kx-3=0\). 要有两个不同的交点，首先判别式 \(\Delta=(-2\sqrt2k)^2-4\left(\frac13-k^2\right)(-3)=4(1-k^2)>0\)，故 \(|k|<1\)；同时二次项系数不能为零，否则只有一个有限交点，所以 \(k\ne\pm\frac{\sqrt3}{3}\). 设两交点横坐标为 \(x_1\)，\(x_2\)，由韦达定理得 \(x_1+x_2=\frac{2\sqrt2k}{\frac13-k^2}\)，\(x_1x_2=-\frac3{\frac13-k^2}\). 由于 \(y_i=kx_i+\sqrt2\)，有 \(\overrightarrow{OA}\cdot\overrightarrow{OB}=x_1x_2+y_1y_2=(1+k^2)x_1x_2+\sqrt2k(x_1+x_2)+2\). 代入上式化简，得到 \(\overrightarrow{OA}\cdot\overrightarrow{OB}-2=\frac{k^2-3}{\frac13-k^2}\). 在 \(|k|<1\) 时分子恒为负，因此该式大于零当且仅当 \(\frac13-k^2<0\)，即 \(|k|>\frac{\sqrt3}{3}\). 故 \(k\in\left(-1,-\frac{\sqrt3}{3}\right)\cup\left(\frac{\sqrt3}{3},1\right)\).
\end{enumerate}
\end{solution}
\end{problem}

\begin{problem}
数列 \(\{a_{n}\}\) 满足 \(a_1 = 1\) 且 \(8a_{n + 1}a_n - 16a_{n + 1} + 2a_n + 5 = 0\)（\(n \geqslant 1\)）. 记 \(b_{n} = \frac{1}{a_{n} - \frac{1}{2}}\)（\(n \geqslant 1\)）.

\begin{enumerate}
\item 求 \(b_{1}\)，\(b_{2}\)，\(b_{3}\)，\(b_{4}\) 的值；
\item 求数列 \(\{b_n\}\) 的通项公式及数列 \(\{a_n b_n\}\) 的前 \(n\) 项和 \(S_n\).
\end{enumerate}

\begin{solution}
\begin{enumerate}
\item 由 \(a_1=1\) 得 \(b_1=\frac1{1-\frac12}=2\). 由递推式解出 \(a_{n+1}=\frac{2a_n+5}{16-8a_n}\)，从而 \(a_{n+1}-\frac12=\frac{3(2a_n-1)}{8(2-a_n)}\). 由 \(a_1\neq\frac12\)，且上式表明只要 \(a_n\neq\frac12\)，\(2\) 就有 \(a_{n+1}\neq\frac12\)；而 \(a_n=2\) 会使原递推式变为 \(9=0\)，不可能，所以归纳可知各项 \(a_n\neq\frac12\). 因而代入 \(b_n=\frac1{a_n-\frac12}\) 可得 \(b_{n+1}=2b_n-\frac43\). 于是 \(b_2=2\times2-\frac43=\frac83\)，\(b_3=2\times\frac83-\frac43=4\)，\(b_4=2\times4-\frac43=\frac{20}{3}\).
\item 令 \(c_n=b_n-\frac43\)，则 \(c_{n+1}=2c_n\)，且 \(c_1=\frac23\)，所以 \(c_n=\frac23\,2^{n-1}=\frac{2^n}{3}\)，从而 \(b_n=\frac{2^n+4}{3}\). 又由 \(b_n=\frac1{a_n-\frac12}\) 得 \(a_n=\frac12+\frac1{b_n}\)，故 \(a_nb_n=\frac{b_n}{2}+1=\frac{2^n+10}{6}\). 因此 \(S_n=\sum_{i=1}^n a_ib_i=\frac16\left(\sum_{i=1}^n2^i+10n\right)=\frac16\left(2^{n+1}-2+10n\right)=\frac{2^n+5n-1}{3}\).
\end{enumerate}
\end{solution}
\end{problem}
